% Drawing 69/05-13 -- folder 69, batch 05, archive page 100 (PDF page 101).
% Transcription: Un petit diagramme : $G$ au centre, avec une flèche vers le haut vers $G/N_1$, une vers la droite vers $G/N_2$, et un trait vers le bas à gauche vers $G/N_3$.
% Kind: diagram G with three quotients. Method 2 (schematic redraw from the facsimile). Status: ok.
% Notes: G in the centre with an arrow up to G/N_1, an arrow to the right to G/N_2, and a stroke down-left to G/N_3.
% Coordinates: 1 unit = 1 mm = 1/100 of the crop width (bbox in page fractions: [0.08, 0.15, 0.4, 0.29]).
% Redrawn by Claude (claude-opus-5-5), issue #73, 2026-11-07; not checked by a human; not compiled (no TeX in the session).
% Rights: reproduces part of the work -- behind RENDER_TRANSCRIPTS, never published by itself.
\documentclass[tikz,border=3pt]{standalone}
\usepackage{amsmath,amssymb}
\usetikzlibrary{patterns,arrows.meta}
\tikzset{pencil/.style={black!45}, struck/.style={black!35}, mark/.style={circle,draw,inner sep=1.3pt}, pt/.style={circle,fill,inner sep=0.9pt}, opt/.style={circle,draw,fill=white,inner sep=1.1pt}}
\begin{document}
\begin{tikzpicture}[x=1mm,y=1mm,line cap=round,line join=round,font=\scriptsize]
\node[inner sep=1pt] at (62,-38) {\scriptsize $G$};
\draw[-{Stealth[length=2mm]}] (62,-32) -- (62,-25);
\draw[-{Stealth[length=2mm]}] (67,-40) -- (79,-45);
\draw (58,-42) -- (52,-48);
\node[inner sep=1pt] at (64,-18) {\scriptsize $G/N_1$};
\node[inner sep=1pt] at (88,-45) {\scriptsize $G/N_2$};
\node[inner sep=1pt] at (44,-51) {\scriptsize $G/N_3$};
\end{tikzpicture}
\end{document}
